\subsection{The quartic gates and the four noncommutative frontiers}
\label{subsec:frontier-nc-quartic-and-tests}

\begin{proposition}
Every complex fusion categorification of any of
$R8\_3\_721$, $R8\_3\_928$, $R8\_3\_941$, and $R8\_3\_1154$
admits a spherical structure and has a pseudo-unitary Galois conjugate.
\end{proposition}
\begin{proof}
The quadratic cases are proved in
Section~\ref{subsec:frontier-nc-galois-pu}.
The same unrestricted signed construction gives the following complete
domains in the two quartic cases:
\[
\begin{array}{c|r|r|r}
\text{ring}&\text{ordered signed tensors}&\text{gauge-original}&
 \text{nongauge excluded}\\\hline
R8\_3\_928&32&16&16\\
R8\_3\_941&128&64&64
\end{array}
\]
No commutativity or modulo-four indicator restriction is imposed.
Each nongauge part is one dual-compatible sign-gauge class. Its
spherical dimension has polynomial
$t^4-320t^3+18688t^2-139264t+278528$, whereas the ordinary scalar
formal codegrees have polynomial
$t^4-576t^3+11520t^2-73728t+147456$.
Induction in the sphericalization requires their quotient to be an
algebraic integer. In a common number field its norm is a positive
integer power of $278528/147456=17/9$, a contradiction.

The scalar algebra of each ordinary ring is the irreducible quartic
FP orbit, and the full complex algebra is
$\mathbb C^4\oplus M_2(\mathbb C)$.
An independent check constructs all five central idempotents and verifies
regular ranks $(1,1,1,1,4)$. For the nongauge signed tensors, rational
regular projectors verify a separable quartic scalar algebra and a
four-dimensional matrix algebra with one-dimensional center.
Its Casimir eigenvalue is $8$, hence its formal codegree is $4$.
Thus no dimension character or matrix formal codegree is omitted.
All surviving signed tensors are gauge-original; the spherical-section
and Galois-transport argument of the preceding proposition applies.
\end{proof}

The canonical quartic replay is
\texttt{noncommutative\_followup/verify\_scope\_gates.py}.
Together the four gates make the following pseudo-unitary calculations
necessary for general complex existence after Galois transport; they
make no assertion of unitarity or existence.

Write $B=MM^{\mathsf T}$ for the induction Gram matrix and the central
restriction matrix. Every restriction column is orthogonal to
$\ker B$. Imposing this exact central-image condition before enumeration
leaves respectively $13,1,1,59$ weighted distinguished frames for
$721,928,941,1154$; ordinary Gram and moment bounds leave $1,1,1,4$.
The two-dimensional representation contributes one distinguished column
with unit entry $2$, whose outer square is counted once.

For an irreducible trace character $\chi_E$, put $m_E=\dim E$.
The distinguished central simple $A_E$ has twist $1$, and its complete
unnormalized $S$ row is
\[
 S_{A_E,Z}=\frac{D}{f_E m_E}\chi_E(FZ),\qquad D=\operatorname{FPdim}(\mathcal C).
\]
The matrix-block character is included. Symmetry and
$STS=DT^{-1}ST^{-1}$ impose
\[
 \sum_Z\chi_E(FZ)\chi_F(FZ)\theta_Z
 =f_Fm_F\chi_E(FA_F).
\]
These supplement the linear indicator equations and the joint
divisor-power rotation constraints. They remove one of two symmetric
assignments for $721$ and one of the last two frames for $1154$.
The surviving $928$ and $941$ branches have exact rational witnesses.
The central eigenvalues $S_{A_E,Z}/\dim Z$ are also tested for algebraic
integrality.

For $721$, complete integral Gram enumeration has four patterns.
Each survives joint rotation through height $12$ and the combined
height-$6$ rotation and $STS$ system with its multiplicities fixed.
For selfdual distinguished simples, the ordinary Verlinde formula gives
\[
 N_{A_E,A_F}^{A_G}
 =D^{-2}\sum_Z
 \frac{S_{A_E,Z}S_{A_F,Z}S_{A_G,Z}}{\dim Z}.
\]
All twenty coefficients are nonnegative integers for every pattern;
the corresponding forgetful tensor equations remain exactly feasible.
For $1154$, the complete integral Gram domain has $14211$ patterns.
Cubic Verlinde integrality leaves $1238$; all pass the ten exact rational
forgetful tensor cones. The subsequent fixed-pattern rotation calculation
is given separately below.

The exact pointed-bimodule stabilizer test also admits both $C_2$
associator classes wherever applicable. The forced pseudo-unitary
invertible indicators are $-1,+1,-1$ for $928,941,1154$, respectively.
These checks do not construct the required module coherences or
associators. Four higher-height numerical attempts remain explicitly
undecided and are not counted as completed obstructions. Exact feasible
relaxations, including those above, are not categorical models.
